% ===========================================================================
\section{Block codes and Shannon's source coding theorem}\label{sec:block-codes}
\warmth{0}\quad\whisper{Encode $n$ symbols at once, and the rounding loss is shared by all of them.}

\begin{strand}
Applying a Shannon code to blocks of $n$ i.i.d.\ symbols keeps the entropy floor
but divides the $+1$ rounding penalty by $n$. Letting $n$ grow shows that the
entropy rate is achievable to any accuracy: this is Shannon's 1948 source coding
theorem.
\end{strand}

\trigger{Use blocks whenever a per-symbol bound is off by an additive constant that does not shrink on its own.}

\block{The basic idea}

Consider the $\X^n$-valued random variable $(X_1,\ldots,X_n)$, where
$X_1,\ldots,X_n$ are i.i.d.\ copies of $X$. Independence makes entropy additive,
\[
  \Ent(X_1,\ldots,X_n)=\answerin[1.6cm]{n\Ent(X)},
\]
\studentrecall{Entropy of an i.i.d. block is $n\Ent(X)$.}
so Proposition~\ref{prop:shannon-code}, applied to the block variable rather
than to a single symbol, gives
\[
  n\Ent(X)\le\E\bigl[|C(X_1,\ldots,X_n)|\bigr]<n\Ent(X)+1
\]
whenever $C$ is an optimal binary code for the block. Dividing by $n$ turns this
into a per-symbol statement:
\[
  \Ent(X)
  \le\frac1n\E\bigl[|C(X_1,\ldots,X_n)|\bigr]
  <\Ent(X)+\answerin[1.2cm]{\tfrac1n}.
\]
\studentrecall{The per-symbol excess is at most $1/n$.}
The excess is now under our control: choose $n$ large and it disappears.

\block{Shannon's theorem}

\begin{theorem}[source coding theorem, Shannon 1948]\label{thm:shannon-source-coding}
Let $X_1,X_2,\ldots$ be an i.i.d.\ sequence of discrete random variables with
state space $\X$. For every $\epsilon>0$ there exist $n\in\mathbb{N}$ and a map
\[
  C:\X^n\longrightarrow\{0,1\}^*
\]
such that:
\begin{enumerate}
  \item the induced map
  $C^*:\bigcup_{k=0}^\infty\X^{nk}\to\{0,1\}^*$ given by
  \[
    (x_1,\ldots,x_{nk})
    \longmapsto
    C(x_1,\ldots,x_n)\,C(x_{n+1},\ldots,x_{2n})
    \cdots C(x_{(k-1)n+1},\ldots,x_{nk})
  \]
  is injective;
  \item $\displaystyle
  \frac1n\E\bigl[|C(X_1,\ldots,X_n)|\bigr]\le\Ent(X)+\epsilon$.
\end{enumerate}
\end{theorem}

\begin{proof}
Choose $n$ with $1/n\le\epsilon$, and regard $Y=(X_1,\ldots,X_n)$ as a single
random variable on the finite alphabet $\X^n$. By
Proposition~\ref{prop:shannon-code} with $D=2$ there is a binary prefix code
$C:\X^n\to\{0,1\}^*$ with
\[
  \Ent(Y)\le\E\bigl[|C(Y)|\bigr]<\Ent(Y)+1 .
\]
Since the $X_i$ are independent, $\Ent(Y)=n\Ent(X)$, so
\[
  \frac1n\E\bigl[|C(Y)|\bigr]<\Ent(X)+\frac1n\le\Ent(X)+\epsilon,
\]
which is the second claim. For the first, note that $C$ is a prefix code, hence
uniquely decodable, so its extension to concatenated blocks is injective: a
decoder reads the encoded string from the left, recovers one block codeword at a
time, and thereby recovers $(x_1,\ldots,x_{nk})$. Thus $C^*$ is injective.
\end{proof}

\begin{yourturn}
Which property of $C$ makes the induced map $C^*$ injective, and why is
nonsingularity of $C$ alone not enough?
\studentrecall{Prefix-freeness; a merely nonsingular code can have two block sequences with the same concatenation.}
\ifshowanswers\else\workspace[2]\fi
\answerprose{Prefix-freeness. It lets the decoder parse the stream block by
block without lookahead. A nonsingular code only distinguishes single blocks;
after concatenation two different block sequences can produce the same string,
as in the code $\{0,01,10\}$.}
\end{yourturn}

\block{Two routes to the same rate}

\noindent\begin{tabularx}{\linewidth}{@{}l L@{}}
\toprule
{\color{indigo}route} & \multicolumn{1}{l}{\color{madder}how entropy appears}\\
\midrule
Shannon code on blocks & lengths $\lceil-\log_2p(x^n)\rceil$ average to
  $n\Ent(X)+O(1)$\\
typical sets (\S\ref{sec:aep}) & index the $\approx2^{n\Ent(X)}$ typical
  sequences with $\approx n\Ent(X)$ bits, and spend a flag bit on the rest\\
\bottomrule
\end{tabularx}

\begin{yourturn}
Using the AEP route, roughly how many bits per symbol are needed to index the
typical set, and what happens to the atypical sequences?
\studentrecall{About $\Ent(X)$ bits per symbol; the atypical set has probability $<\epsilon$ and can be encoded wastefully.}
\ifshowanswers\else\workspace[2]\fi
\answerprose{Indexing $|T_{n,\epsilon}|\le2^{n(\Ent(X)+\epsilon)}$ sequences needs
about $\Ent(X)+\epsilon$ bits per symbol. The atypical sequences carry
probability less than $\epsilon$, so even a wasteful $\log_2|\X|$ bits per symbol
for them contributes a vanishing amount to the average.}
\end{yourturn}

\block{Beyond i.i.d.: the entropy rate}

Everything above used independence only once, to write
$\Ent(X_1,\ldots,X_n)=n\Ent(X)$. For a dependent source that quantity is
replaced by a limit.

\begin{definition}[entropy rate]\label{def:entropy-rate}
For a stationary sequence $X_1,X_2,\ldots$ the \keyword{entropy rate} is
\[
  \bar\Ent=\lim_{n\to\infty}\answerin[2.4cm]{\tfrac1n\Ent(X_1,\ldots,X_n)},
\]
which exists because $\Ent(X_1,\ldots,X_n)$ is subadditive in $n$.
\studentrecall{Per-symbol entropy of a block, in the limit: $\frac1n\Ent(X_1,\dots,X_n)$.}
\end{definition}

\begin{theorem}[Shannon--McMillan--Breiman]\label{thm:smb}
Let $X_1,X_2,\ldots$ be an ergodic stationary sequence on a finite state space
$\X$. Then
\[
  -\frac1n\log p_{X^n}(X_1,\ldots,X_n)
  \longrightarrow\bar\Ent
  \qquad\text{in probability, as }n\to\infty,
\]
where $\bar\Ent=\lim_{n\to\infty}\frac1n\Ent(X_1,\ldots,X_n)$.
\end{theorem}

This is the AEP with independence weakened to \emph{ergodic stationarity} and
$\Ent(X)$ replaced by $\bar\Ent$; the statement even holds almost surely. Once
it is available, typical sets, block coding and
Theorem~\ref{thm:shannon-source-coding} all go through verbatim with the
entropy rate in place of the single-symbol entropy.

\begin{yourturn}
For an i.i.d.\ sequence, what does $\bar\Ent$ reduce to, and which hypothesis of
Theorem~\ref{thm:smb} replaces independence?
\studentrecall{$\bar\Ent=\Ent(X_1)$; independence is weakened to ergodic stationarity.}
\ifshowanswers\else\workspace[2]\fi
\answerprose{Additivity gives $\Ent(X_1,\ldots,X_n)=n\Ent(X_1)$, so
$\bar\Ent=\Ent(X_1)$ and the theorem becomes the AEP. Independence is replaced
by the assumption that the sequence is stationary and ergodic, which is what
lets the time average of the surprise converge to its ensemble mean.}
\end{yourturn}

\vspace{0.8cm}
\recall{Why does blocking remove the $+1$ penalty but never the entropy floor?}
